\subsection{Quantum Class 5 model and spatial Lax matrix}
\label{sec:class5-spatial-lax}

The name ``Class 5'' refers to one of the regular $4\times4$ difference-form
solutions of the Yang--Baxter equation found in the classification of
nearest-neighbour integrable spin-$\tfrac12$ chains in
Ref.~\cite{deLeeuwPribytokRyan2019}.  Its continuum limit was subsequently
studied in Ref.~\cite{deLeeuwFontanellaNietoGarcia2026}, where it was shown to
be a non-Hermitian integrable deformation of the XXX spin chain whose
long-wavelength limit is a deformed Landau--Lifshitz model.  We use a
one-parameter representative of this model for the quantum-to-classical
Lax calculation.  This representative is
related to the general Class 5 $R$-matrix by the spectral normalization and
local basis transformation described in
Ref.~\cite{deLeeuwFontanellaNietoGarcia2026}.

Let $\kappa_5\in\mathbb C$ denote the \emph{Class 5 lattice deformation parameter}.
It corresponds to the parameter denoted by $a$ in
Ref.~\cite{deLeeuwFontanellaNietoGarcia2026}.  In the
ordered two-spin basis
\begin{equation}
\bigl\{
|\!\uparrow\uparrow\rangle,
|\!\uparrow\downarrow\rangle,
|\!\downarrow\uparrow\rangle,
|\!\downarrow\downarrow\rangle
\bigr\},
\label{eq:class5-two-spin-basis}
\end{equation}
the two-site Hamiltonian operator defining the Class 5 chain in our
normalization is
\begin{equation}
h_5(\kappa_5)
=
\begin{pmatrix}
2&2\kappa_5&-2\kappa_5&0\\
0&0&2&0\\
0&2&0&0\\
0&0&0&2
\end{pmatrix}.
\label{eq:class5-local-hamiltonian}
\end{equation}
Following the embedding convention of Section~\ref{sec:quantum-lattice-setup},
the copy of $h_5$ on the bond joining sites $n$ and $n+1$ is denoted by
$h_{5;n,n+1}(\kappa_5)$.  The periodic quantum Hamiltonian is therefore
\begin{equation}
H^{\rm lat}_5(\kappa_5)
=
\sum_{n=1}^{N}h_{5;n,n+1}(\kappa_5),
\qquad
h_{5;N,N+1}\equiv h_{5;N,1}.
\label{eq:class5-lattice-hamiltonian}
\end{equation}
For $\kappa_5\neq0$ this Hamiltonian is non-Hermitian in the standard inner product.
The undeformed value $\kappa_5=0$ reduces, up to the normalization fixed in
Section~\ref{sec:quantum-lattice-setup}, to the spin-$\tfrac12$ XXX chain.

The same normalization is generated by the regular Class 5 quantum
$R$-matrix
\begin{equation}
R_5(u;\kappa_5)
=
\begin{pmatrix}
1+2u&2\kappa_5 u&-2\kappa_5 u&0\\
0&2u&1&0\\
0&1&2u&0\\
0&0&0&1+2u
\end{pmatrix},
\label{eq:class5-R-matrix}
\end{equation}
where $u$ is the quantum spectral parameter introduced in
Section~\ref{sec:quantum-lattice-setup}.  It obeys the Yang--Baxter equation
and the regularity condition
\begin{equation}
R_5(0;\kappa_5)=P,
\label{eq:class5-R-regularity}
\end{equation}
and a direct differentiation gives
\begin{equation}
h_5(\kappa_5)
=
P\left.\frac{\partial R_5(u;\kappa_5)}{\partial u}\right|_{u=0},
\label{eq:class5-h-from-R}
\end{equation}
in agreement with the general relation
\eqref{eq:h-from-R}.  Equations~\eqref{eq:class5-local-hamiltonian} and
\eqref{eq:class5-R-matrix} fix the model and the normalizations used in this
section.

We isolate the deformation in the two-site operator $K_5$, defined by
\begin{equation}
K_5
:=
\begin{pmatrix}
0&1&-1&0\\
0&0&0&0\\
0&0&0&0\\
0&0&0&0
\end{pmatrix}
\in\End(\mathcal H_{\rm loc}\otimes\mathcal H_{\rm loc}).
\label{eq:class5-K-definition}
\end{equation}
With the permutation operator $P$ defined in
\eqref{eq:local-permutation-operator}, the Class 5 data become
\begin{equation}
R_5(u;\kappa_5)=2u\Id_4+P+2\kappa_5 uK_5,
\qquad
h_5(\kappa_5)=2P+2\kappa_5 K_5.
\label{eq:class5-R-h-PK}
\end{equation}
The operator $K_5$ is nilpotent and satisfies
\begin{equation}
K_5^2=0,
\qquad
PK_5=K_5,
\qquad
K_5P=-K_5.
\label{eq:class5-PK-algebra}
\end{equation}
These identities simplify the finite-lattice time Lax operator.

We now specify the long-wavelength scaling.  The lattice deformation is sent
to zero together with the lattice spacing $\epsilon$, while the
\emph{Class 5 continuum deformation coupling} $\alpha_5$ is held fixed:
\begin{equation}
\kappa_5=\epsilon\alpha_5.
\label{eq:class5-deformation-scaling}
\end{equation}
This is the scaling used in the direct continuum construction of
Ref.~\cite{deLeeuwFontanellaNietoGarcia2026}; it places the deformation and
the ordinary Landau--Lifshitz terms at the same order in the continuum
action.  For the Lax operator we simultaneously relate the quantum spectral
parameter $u$ to the continuum Lax spectral parameter $\lambda$ by
\begin{equation}
u=\frac{\lambda}{\epsilon},
\qquad \lambda\neq0,
\label{eq:class5-spectral-scaling}
\end{equation}
and use the common time scaling $t=\epsilon^2t_{\rm lat}$ introduced in
\eqref{eq:common-time-scaling}.

For the fundamental spin-$\tfrac12$ representation the lattice Lax operator
is $L_{5;a,n}(u):=R_{5;a,n}(u;\kappa_5)$.  We normalize it by its leading scalar
part,
\begin{equation}
\widehat L_{5;a,n}(\epsilon,\lambda)
:=
\frac{1}{2u}L_{5;a,n}(u).
\label{eq:class5-normalized-L}
\end{equation}
The quantum spectral derivatives are taken at fixed $\kappa_5$, before
applying the scalings \eqref{eq:class5-deformation-scaling} and
\eqref{eq:class5-spectral-scaling}.  The normalized operator has the exact
expansion
\begin{equation}
\widehat L_{5;a,n}
=
\Id_{\mathcal V_a\otimes\mathcal H_n}
+\epsilon\,L^{(1)}_{5;a,n},
\qquad
L^{(1)}_5
=
\frac{P}{2\lambda}+\alpha_5 K_5.
\label{eq:class5-L-expansion}
\end{equation}
Here $L^{(1)}_{5;a,n}$ denotes the copy of the fixed two-site operator $L^{(1)}_5$
on $\mathcal V_a\otimes\mathcal H_n$.

To evaluate its lower symbol, define the classical transverse spin
combinations
\begin{equation}
S^+(x,t):=S^1(x,t)+\ii S^2(x,t),
\qquad
S^-(x,t):=S^1(x,t)-\ii S^2(x,t).
\label{eq:class5-spin-pm}
\end{equation}
Throughout this section, $S^{\pm}$ denotes these components of the classical
unit spin field defined in \eqref{eq:continuum-spin-field}.  At the point $x=x_n$, the lower symbols of
$P_{a,n}$ and $K_{5;a,n}$ are
\begin{align}
P_{a,n}^{\downarrow}
&=
\frac12
\begin{pmatrix}
1+S^3&S^-\\
S^+&1-S^3
\end{pmatrix},
\label{eq:class5-P-lower-symbol}\\
K_{5;a,n}^{\downarrow}
&=
\frac12
\begin{pmatrix}
S^+&-(1+S^3)\\
0&0
\end{pmatrix}.
\label{eq:class5-K-lower-symbol}
\end{align}
Substituting these expressions into \eqref{eq:class5-L-expansion} gives
\begin{equation}
\widehat L_{5;a,n}^{\downarrow}
=
\Id_{\mathcal V_a}
+\epsilon U_5(x_n,t;\lambda),
\label{eq:class5-L-lower-expansion}
\end{equation}
where the \emph{Class 5 continuum spatial Lax matrix} is
\begin{equation}
U_5(x,t;\lambda)
=
\frac{1}{4\lambda}
\begin{pmatrix}
1+S^3+2\alpha_5\lambda S^+
&
S^- -2\alpha_5\lambda(1+S^3)\\
S^+
&
1-S^3
\end{pmatrix}.
\label{eq:class5-spatial-Lax}
\end{equation}
The finite-lattice limit retains the complete $\mathfrak{gl}(2)$ matrix in
Eq.~\eqref{eq:class5-spatial-Lax}, including its scalar part.
